\documentclass[11pt,a4paper]{article}
\usepackage[T1]{fontenc}
\usepackage{isabelle,isabellesym}
\usepackage{amsmath,amssymb}
\usepackage{pdfsetup}

\urlstyle{rm}
\isabellestyle{it}

\begin{document}

\title{Pascal's Hexagon Theorem}
\author{Ari Gilchrist}
\maketitle

\begin{abstract}
  We formalise Pascal's hexagon theorem (1640): for any six distinct points on a
  non-degenerate conic, the three intersection points of the opposite sides of the hexagon
  they form are collinear, provided those intersections exist.
  The main result is stated in the real plane $\mathbb{R}^2$ using the standard notion of
  collinearity from HOL-Analysis. We also prove several variants of the theorem,
  including the case of arbitrary and possibly degenerate conics, and the special cases in which
  the conic is a pair of lines (Pappus's hexagon theorem) or a circle.
\end{abstract}

\subsection*{Acknowledgements}
This formalisation was generated almost entirely by Anthropic's Claude
Opus 5.5, which designed the proof strategy, wrote and debugged the Isabelle code, and
drafted the accompanying document. The author only reviewed the development,
revised the text, and verified that the definitions and theorem statements capture the intended
mathematics.

\tableofcontents

\parindent 0pt\parskip 0.5ex

\section{Introduction}

Let $A, B, C, D, E, F$ be six arbitrary distinct points on a conic, defining the
hexagon $ABCDEF$. Its \emph{opposite sides} are the pairs $AB, DE$; $BC, EF$ and $CD, FA$.
The vertices may lie on the conic in any order, hence the hexagon may cross over itself.
Pascal~\cite{pascal1640} observed that the three intersection points of opposite
sides always lie on a common line, now called the \emph{Pascal line} of the hexagon:

\begin{quote}
  \textbf{Theorem (Pascal, 1640).} Let $A, \dots, F$ be six distinct points on a
  non-degenerate conic in $\mathbb{R}^2$. If $P$ lies on the lines $AB$ and $DE$, $Q$ on
  $BC$ and $EF$, and $R$ on $CD$ and $FA$, then $P$, $Q$ and $R$ are collinear.
\end{quote}

Here a non-degenerate conic is the zero set of a polynomial
$ax^2 + bxy + cy^2 + dx + ey + f$ whose homogenisation has an invertible symmetric
matrix; if it is non-empty, it is an ellipse, a parabola or a hyperbola. The theorem
assumes that $P$, $Q$ and $R$ exist, i.e.\ that no two opposite sides are parallel.

We follow the standard coordinate proof in the projective plane (see
e.g.~\cite{richter-gebert2011} for background on homogeneous coordinates and projective
conics). A point is represented by a homogeneous coordinate vector, the line through two
points is their cross product, and three points are collinear iff the determinant of their
coordinate vectors vanishes. The collinearity of the three intersection points is
therefore expressed by a single \emph{Pascal determinant}. A linear change of coordinates
moves $A$, $B$, $C$ to the standard basis vectors, which multiplies the Pascal determinant
by a non-zero factor and turns the conic into $\alpha xy + \beta yz + \gamma zx = 0$. The
three points $D$, $E$, $F$ on the conic then give three linear equations satisfied by the
non-zero vector $(\alpha, \beta, \gamma)$, so the matrix of this system is singular; in
these coordinates its determinant is, up to sign, the Pascal determinant. This argument works
over any field and for any conic, provided that $A$, $B$, $C$ are not collinear and that
each pair of opposite sides consists of two distinct lines; both are ensured by seven
non-collinearity conditions on the vertices.

Section~2 develops this projective argument in homogeneous coordinates over an arbitrary
field. It also proves that a line meets a non-degenerate conic in at most two points, so
that for such conics the non-collinearity conditions follow from distinctness of the six
points. Section~3 transfers the results to $\mathbb{R}^2$ via $(x, y) \mapsto (x, y, 1)$,
using the standard notions \texttt{collinear} and \texttt{dist} from HOL-Analysis. Besides
Pascal's theorem as stated above (\texttt{pascal\_hexagon}), it proves:
\begin{itemize}
\item \texttt{pascal\_hexagon\_any\_conic}: the theorem for an arbitrary conic, possibly
  degenerate, assuming that no three of the six points are collinear except possibly the
  alternate triples $A, C, E$ and $B, D, F$. Up to renaming the points and reordering the
  coefficients of the conic, this is the statement of \texttt{PASCAL\_AFFINE} in
  HOL Light~\cite{harrison-pascal}.
\item \texttt{pascal\_hexagon\_circle}: the theorem for six distinct points on a circle.
  Up to renaming, this is the statement of \texttt{PASCAL\_AFFINE\_CIRCLE} in HOL Light and
  of the Lean formalisation~\cite{alexeev-pascal}.
\item \texttt{pappus\_hexagon}: Pappus's hexagon theorem, the case in which the conic is
  a pair of lines with $A, C, E$ on one and $B, D, F$ on the other.
\end{itemize}

Pascal's theorem is number 28 on Freek Wiedijk's list ``Formalizing 100
Theorems''~\cite{wiedijk100}. It has previously been formalised in HOL Light by John
Harrison~\cite{harrison-pascal}, in Mizar by Roland Coghetto~\cite{coghetto2017}, in Rocq
by Nicolas Magaud and Julien Narboux~\cite{magaud-narboux-pascal}, and in Lean by Boris
Alexeev~\cite{alexeev-pascal}. In Isabelle/HOL, Bordg's synthetic development of projective
geometry~\cite{Projective_Geometry-AFP} proves the analogue of the Rocq result: a projective
plane has Pappus's property if and only if, whenever the opposite sides of a hexagon meet in
collinear points, so do those of the hexagon obtained by swapping two adjacent vertices.
IsaGeoCoq~\cite{IsaGeoCoq-AFP} proves the parallel form of Pappus's theorem in Tarski's
axiomatisation of Euclidean geometry. Neither development involves conics, and to the
author's knowledge, Pascal's theorem for conics has not previously been formalised in
Isabelle/HOL.

\input{session}

\bibliographystyle{abbrv}
\bibliography{root}

\end{document}
